% ---------------------------------------------------------------------------
% AMR dans un atelier de câblage -- Planification, algorithmes, méthodes,
% bibliothèques et justification des choix.
%
% Compilation : pdflatex (2 passes, pour la table des matières), par ex. sur
% Overleaf : téléverser ce fichier + les images de ..\media\ dans un dossier
% "figures/" (ou garder l'arborescence du projet). Les figures absentes sont
% remplacées par un cadre, le document compile quand même.
% ---------------------------------------------------------------------------
\documentclass[11pt,a4paper]{article}

\usepackage[utf8]{inputenc}
\usepackage[T1]{fontenc}
\usepackage[french]{babel}
\usepackage{lmodern}
\usepackage[margin=2.3cm]{geometry}
\usepackage{amsmath,amssymb}
\usepackage{graphicx}
\usepackage{booktabs}
\usepackage{tabularx}
\usepackage{array}
\usepackage{xcolor}
\usepackage{enumitem}
\usepackage{algorithm}
\usepackage{algpseudocode}
\usepackage{tikz}
\usetikzlibrary{arrows.meta,positioning,shapes.geometric,fit}
\usepackage[hidelinks]{hyperref}

\graphicspath{{../media/}{figures/}{media/}}
\setlist{itemsep=2pt,topsep=4pt}
\renewcommand{\arraystretch}{1.25}
\floatname{algorithm}{Algorithme}
\renewcommand{\algorithmicrequire}{\textbf{Entrées :}}
\renewcommand{\algorithmicensure}{\textbf{Sortie :}}

% figure facultative : image si elle existe, sinon un cadre avec son nom
\newcommand{\figimg}[2]{%
  \IfFileExists{../media/#2}{\includegraphics[width=#1\linewidth]{#2}}{%
  \IfFileExists{figures/#2}{\includegraphics[width=#1\linewidth]{#2}}{%
  \IfFileExists{media/#2}{\includegraphics[width=#1\linewidth]{#2}}{%
  \fbox{\parbox{#1\linewidth}{\centering\vspace{1.5cm}\texttt{#2}\\(image à ajouter)\vspace{1.5cm}}}}}}}

\newcommand{\code}[1]{\texttt{#1}}
\newcolumntype{L}{>{\raggedright\arraybackslash}X}

\title{\textbf{Robot mobile autonome (AMR) dans un atelier de câblage}\\[4pt]
\large Planification du projet, algorithmes et méthodes,\\ bibliothèques utilisées et justification des choix}
\author{Nassim Chawad\\ Module « Atelier intelligent »}
\date{Novembre 2026}

\begin{document}
\maketitle

\begin{abstract}
Ce document décrit comment a été conçu et construit un robot mobile autonome (AMR) qui transporte
des bacs de matière (bobines, faisceaux) entre les postes d'un atelier de câblage simulé dans
CoppeliaSim. Il présente la démarche suivie (découpage en parties, tests à chaque étape), puis,
pour chaque fonction du robot (localisation, carte, planification de chemin, suivi du chemin et
évitement d'obstacles, suivi et prédiction du mouvement des opérateurs, accostage, manipulation,
cellule de contrôle qualité), l'algorithme retenu, sa
formulation, ses paramètres et la raison de ce choix face aux alternatives. Une dernière partie
résume les résultats des tests et les limites connues.
\end{abstract}

\tableofcontents
\newpage

% ===========================================================================
\section{Contexte et cahier des charges}

L'objectif est de simuler un AMR dans une usine de faisceaux électriques : le robot reçoit des
ordres de transport (« bobines : Magasin $\rightarrow$ Coupe »), \emph{planifie lui-même} son
chemin, \emph{évite} les obstacles imprévus, notamment des opérateurs qui marchent, s'accoste
précisément aux postes, charge et décharge les bacs avec son bras, et diffuse un flux caméra en
temps réel. Aucune contrainte technique n'était imposée : tous les choix ci-dessous sont des
décisions de conception, prises au début du projet (partie 0) puis affinées.

\paragraph{Décisions structurantes (partie 0).}
\begin{enumerate}[label=\textbf{\Alph*.}]
  \item \textbf{Carte connue, pas de SLAM.} Le plan de l'atelier est fixe et connu : la carte est
        générée à partir de la description de l'atelier. Le lidar sert à détecter ce qui \emph{n'est
        pas} sur la carte (personnes, objets égarés).
  \item \textbf{Trois capteurs obligatoires :} lidar 2D (perception, accostage), ceinture
        d'ultrasons (sécurité à courte portée, tout autour du robot), caméra (flux vidéo,
        identification des postes par marqueurs ArUco).
  \item \textbf{Chemin bloqué :} ralentir, s'arrêter et attendre, puis re-planifier après
        quelques secondes ; s'écarter si une personne continue d'approcher.
  \item \textbf{Scène lisible :} six postes de production + un dock de recharge, des rayonnages, des
        opérateurs ; la scène doit clairement ressembler à une usine de câblage.
  \item \textbf{Supervision :} à terme, une interface de suivi des ordres de transport et
        des indicateurs d'évaluation.
\end{enumerate}

\paragraph{Flux de production simulé.}
Magasin (bobines brutes) $\rightarrow$ Coupe/dénudage $\rightarrow$ Sertissage $\rightarrow$
Assemblage (planches de câblage) $\rightarrow$ Test électrique (cellule de contrôle qualité)
$\rightarrow$ Expédition, plus une table \textsc{Rebuts} (retour en Assemblage pour reprise) et le
dock de recharge de l'AMR. L'atelier mesure $20\times12$\,m, avec une allée centrale de 2,5\,m.

\begin{figure}[h]
  \centering
  \figimg{0.85}{atelier_vue_3d.png}
  \caption{L'atelier de câblage simulé (CoppeliaSim).}
\end{figure}

% ===========================================================================
\section{Démarche de planification du projet}

\subsection{Découpage en parties testées une à une}
Le projet a été construit \emph{incrémentalement} : chaque partie ajoute une seule capacité, avec
un programme de test qui la valide avant de passer à la suivante. Cela évite d'empiler des erreurs
et rend chaque résultat mesurable.

\begin{table}[h]
\centering\small
\begin{tabularx}{\linewidth}{@{}lLL@{}}
\toprule
\textbf{Partie} & \textbf{Contenu} & \textbf{Validation} \\
\midrule
0  & Cahier des charges, choix A--E, arborescence, journal d'installation & --- \\
1  & Liaison Python $\leftrightarrow$ CoppeliaSim (API ZMQ, mode pas-à-pas) & \code{test\_connection} \\
2  & Pilote du youBot (cinématique mécanum, localisation simulée) & \code{test\_drive}, \code{test\_square} \\
3  & Capteurs : lidar, 8 ultrasons, caméra & \code{test\_sensors} \\
4  & Scène de l'atelier, générée depuis une description unique & images de la scène \\
5  & Carte d'occupation + planificateur A* & \code{test\_planner} (28/28 trajets) \\
6  & Suivi de chemin + évitement des opérateurs ; puis suivi des opérateurs (filtre de Kalman) et
     navigation par prédiction & \code{test\_navigation}, \code{test\_tracking}, \code{test\_errand} \\
7  & Accostage ArUco + lidar, bras, prise et dépose & \code{test\_aruco}, \code{test\_docking}, \code{test\_pick\_place}, \code{test\_errand} \\
8  & Gestionnaire de missions (ordres, batterie) & \emph{à venir} \\
9  & Cellule de contrôle qualité (caméra + UR5) & \code{test\_qc\_cell} \\
10 & Interface de supervision & \emph{à venir} \\
11 & Évaluation, rapport, vidéo & \emph{à venir} \\
\bottomrule
\end{tabularx}
\caption{Plan du projet et test de validation de chaque partie.}
\end{table}

\subsection{Principes de travail}
\begin{itemize}
  \item \textbf{Une seule source de vérité.} Tout l'atelier est décrit comme des données dans
        \code{layout.py} (formes, modèles de bibliothèque, postes, opérateurs). Le même fichier sert
        à construire la scène (\code{build\_scene.py}) \emph{et} à générer la carte du robot
        (\code{occupancy.py}) : la scène et la carte ne peuvent pas diverger.
  \item \textbf{Simulation pas-à-pas.} Python commande le simulateur pas par pas
        ($\Delta t = 50$\,ms) : chaque décision est prise sur un état cohérent, les tests sont
        reproductibles et la vitesse de calcul de Python n'influence pas la physique.
  \item \textbf{Mesurer plutôt que supposer.} Les dimensions des modèles de bibliothèque, les
        paramètres du bras, la hauteur de la plateforme, etc. ont été mesurés dans le simulateur
        (scripts d'inspection) avant d'être utilisés.
  \item \textbf{Traçabilité.} Chaque installation est consignée (version, emplacement, commande de
        désinstallation) ; le code est versionné avec Git ; une note de rapport est écrite à la fin
        de chaque partie.
\end{itemize}

% ===========================================================================
\section{Environnement logiciel et bibliothèques}

\subsection{Choix et justification}
\begin{table}[h]
\centering\small
\begin{tabularx}{\linewidth}{@{}lp{3.2cm}L@{}}
\toprule
\textbf{Outil} & \textbf{Rôle} & \textbf{Pourquoi ce choix} \\
\midrule
CoppeliaSim Edu 4.10 & Le « monde » : scène, physique, robot, capteurs simulés &
  Gratuit pour l'enseignement, bibliothèque riche (youBot, UR5, lidar Hokuyo, convoyeurs, personnages
  animés, marqueurs ArUco), plusieurs moteurs physiques, API externe. \\
Python 3 & Le « cerveau » : perception, planification, commande &
  Débogage simple, écosystème scientifique (NumPy, OpenCV), code lisible pour le rapport. \\
API ZMQ Remote & Liaison Python $\leftrightarrow$ simulateur &
  Client officiel, appelle toute l'API de CoppeliaSim depuis Python, mode pas-à-pas synchrone. \\
NumPy & Calcul vectoriel &
  Traitement des nuages de points lidar, transformations de repères, grilles. \\
OpenCV & Vision et images &
  Détection ArUco + \code{solvePnP}, segmentation couleur HSV, dilatation et transformée de
  distance de la carte, affichage temps réel. \\
simIK (CoppeliaSim) & Cinématique inverse du bras UR5 &
  Solveur numérique intégré (moindres carrés amortis), adapté à un bras à 6 axes. \\
Lua (dans la scène) & Petits scripts embarqués &
  Animation des opérateurs, crochets de commande (pince, ventouse, convoyeur). \\
\bottomrule
\end{tabularx}
\caption{Outils et bibliothèques.}
\end{table}

\subsection{Pourquoi pas ROS ?}
L'architecture est \emph{modulaire et inspirée de ROS} (un module par fonction : capteurs, carte,
planificateur, navigateur, accostage, bras), mais sans intergiciel. Pour un projet individuel en
simulation, ROS~2 aurait ajouté une installation lourde (Linux ou WSL), la gestion des
\emph{topics}, des messages et des \emph{launch files}, sans bénéfice fonctionnel ici : un seul
programme commande un seul robot. Les modules Python correspondent chacun à un nœud ROS classique
(\code{map\_server}, \code{global\_planner}, \code{local\_planner}, etc.) ; un passage à ROS~2 via
le plugin \code{simROS2} de CoppeliaSim reste possible.

\subsection{Pourquoi un planificateur maison plutôt qu'OMPL ?}
CoppeliaSim intègre OMPL (planificateurs par échantillonnage : RRT, PRM\dots). Ces méthodes sont
faites pour les grands espaces de configuration (bras à 6 axes). Pour un robot qui se déplace dans
un plan avec une carte en grille, A* est \emph{optimal}, \emph{complet} et \emph{déterministe} (même
chemin à chaque exécution), rapide (environ 15\,ms par trajet ici), et il prend directement en compte
un coût de proximité des obstacles. Écrit à la main, il est aussi entièrement explicable dans ce
rapport.

\begin{figure}[h]
\centering
\begin{tikzpicture}[node distance=6mm and 9mm, >=Stealth,
  box/.style={draw, rounded corners, fill=blue!6, minimum width=2.6cm, minimum height=8mm, align=center, font=\small},
  sim/.style={box, fill=orange!12}]
  \node[sim] (scene) {CoppeliaSim\\scène + physique};
  \node[box, right=of scene] (sens) {Capteurs\\\code{sensors.py}};
  \node[box, right=of sens] (perc) {Perception\\obstacles imprévus};
  \node[box, below=of perc] (nav) {Navigateur\\machine à états};
  \node[box, below=of sens] (plan) {A*\\\code{planner.py}};
  \node[box, below=of scene] (map) {Carte\\\code{occupancy.py}};
  \node[box, below=of nav] (dock) {Accostage\\ArUco + lidar};
  \node[box, left=of dock] (arm) {Bras youBot\\IK + pince};
  \node[box, left=of arm] (drv) {Pilote youBot\\mécanum + UWB};
  \node[box, above=of scene, xshift=-0mm] (lay) {\code{layout.py}\\(source unique)};
  \draw[->] (lay) -- (scene);
  \draw[->] (lay.west) -| ++(-0.6,0) |- (map.west);
  \draw[->] (scene) -- (sens);
  \draw[->] (sens) -- (perc);
  \draw[->] (perc) -- (nav);
  \draw[->] (map) -- (plan);
  \draw[->] (plan) -- (nav);
  \draw[->] (nav) -- (dock);
  \draw[->] (dock) -- (arm);
  \draw[->] (arm) -- (drv);
  \draw[->] (drv.west) -| ++(-0.5,0) |- (scene.south west);
\end{tikzpicture}
\caption{Architecture logicielle (un module Python par fonction).}
\end{figure}

% ===========================================================================
\section{Robot, cinématique et localisation}

\subsection{Choix du robot : KUKA youBot}
Le youBot (bibliothèque CoppeliaSim) a été préféré au Pioneer P3-DX initialement envisagé :
\begin{itemize}
  \item \textbf{roues mécanum} : déplacement omnidirectionnel (avant, latéral, rotation, ou les
        trois à la fois). C'est très utile pour s'écarter d'une personne sans tourner et pour
        l'accostage fin ;
  \item \textbf{bras à 5 axes + pince} : le robot charge et décharge réellement les bacs ;
  \item il n'a pas de capteurs : on y monte lidar, ultrasons et caméra (partie 3).
\end{itemize}

\subsection{Cinématique mécanum}
Chaque roue porte des galets inclinés à $45^\circ$. Pour une consigne de vitesse dans le repère du
robot $(v_x, v_y, \omega)$ (avant, gauche, rotation), avec $R = 0{,}05$\,m le rayon de roue et
$L_x = 0{,}228$\,m, $L_y = 0{,}158$\,m les demi-entraxes :
\begin{equation}
\begin{aligned}
\omega_{av,g} &= \tfrac{1}{R}\big(v_x - v_y - (L_x+L_y)\,\omega\big), &
\omega_{av,d} &= \tfrac{1}{R}\big(v_x + v_y + (L_x+L_y)\,\omega\big),\\
\omega_{ar,g} &= \tfrac{1}{R}\big(v_x + v_y - (L_x+L_y)\,\omega\big), &
\omega_{ar,d} &= \tfrac{1}{R}\big(v_x - v_y + (L_x+L_y)\,\omega\big).
\end{aligned}
\end{equation}
Validation : 0,4\,m demandés $\rightarrow$ 0,392\,m en avant et 0,388\,m en latéral ; 1,0\,rad
$\rightarrow$ 1,006\,rad.

\subsection{Commande de position}
Pour rejoindre une pose $(x^*, y^*, \theta^*)$, l'erreur de position est exprimée dans le repère du
robot et une loi proportionnelle saturée commande translation et rotation \emph{en même temps} :
\[
\begin{pmatrix} e_f \\ e_l \end{pmatrix} =
\begin{pmatrix} \cos\theta & \sin\theta \\ -\sin\theta & \cos\theta \end{pmatrix}
\begin{pmatrix} x^*-x \\ y^*-y \end{pmatrix},\qquad
v = \min(v_{max},\, k_v\,d),\qquad \omega = \operatorname{sat}\big(k_\omega\,(\theta^*-\theta)\big).
\]
Test : carré de 1\,m parcouru en crabe (cap constant), environ 4\,s par côté, erreur réelle de 2 à 6\,cm.

\begin{figure}[h]
  \centering
  \figimg{0.55}{test_square.png}
  \caption{Carré de 1\,m en déplacement latéral : trajectoire réelle et mesurée.}
\end{figure}

\subsection{Localisation : balises UWB simulées, sans odométrie ni SLAM}
\textbf{Choix :} la position vient d'un système de positionnement intérieur simulé
(balises \emph{UWB} : vraie position + bruit gaussien $\sigma = 2$\,cm) et le cap d'un gyroscope
(bruit $\sigma = 0{,}5^\circ$).

\textbf{Justification :}
\begin{itemize}
  \item l'odométrie des roues mécanum est peu fiable (les galets glissent par construction) et dérive ;
  \item un SLAM complet n'apporte rien quand le plan de l'usine est connu et fixe ; c'est d'ailleurs
        le cas des AMR industriels, qui sont livrés avec la carte du site ;
  \item les systèmes UWB existent réellement en usine (précision de l'ordre de quelques
        centimètres) ;
  \item la précision fine nécessaire au bras est obtenue \emph{localement} à l'accostage
        (ArUco + lidar, section~\ref{sec:dock}), pas par la localisation globale.
\end{itemize}
Évolution possible : localisation par le lidar sur la carte connue (AMCL, filtre particulaire).

% ===========================================================================
\section{Capteurs et perception}

\begin{table}[h]
\centering\small
\begin{tabularx}{\linewidth}{@{}lLL@{}}
\toprule
\textbf{Capteur} & \textbf{Caractéristiques} & \textbf{Usage} \\
\midrule
Lidar Hokuyo URG-04LX (bibliothèque) & $240^\circ$, 684 points, 5\,m, à l'avant, 13\,cm du sol &
  Obstacles imprévus, mesure de la face des tables (accostage) \\
8 ultrasons (copiés du P3-DX) & Ceinture tout autour, cônes de $30^\circ$, portée 1\,m &
  Angles morts du lidar (côtés, arrière), arrêt d'urgence à 0,22\,m \\
Caméra (capteur de vision) & $640\times480$, $60^\circ$, inclinée de $15^\circ$ vers le bas &
  Flux vidéo, marqueurs ArUco (identification + pose des postes) \\
\bottomrule
\end{tabularx}
\caption{Capteurs montés sur le youBot.}
\end{table}

\paragraph{Points lidar.} Chaque mesure (profondeur) est transformée du repère du capteur au repère
du robot par la matrice de pose du capteur, $p_{robot} = R\,p_{capteur} + t$ ; les points tombant
sur le robot lui-même sont éliminés (boîte $0{,}36\times0{,}25$\,m). Validation : boîte attendue à
0,95\,m mesurée à 0,950\,m, poteau attendu à 0,64\,m mesuré à 0,641\,m.

\paragraph{Couverture et angles morts.} Les champs des capteurs ne sont pas recopiés à la main : ils
sont \emph{lus dans la scène} (pose de chaque capteur par rapport au robot, volume de détection des
ultrasons, angle et portée des capteurs de profondeur du lidar) et dessinés dans la vue de dessus
temps réel. La figure~\ref{fig:couverture} montre que l'arrière du robot n'est couvert que par trois
cônes d'ultrasons et qu'il reste des secteurs vides entre les cônes.

\paragraph{Pas de lidar arrière (choix assumé).} Le Hokuyo couvre $240^\circ$ par construction et,
monté à 13\,cm à l'avant, il aurait de toute façon le corps du robot dans son champ vers l'arrière.
Les AMR industriels utilisent deux lidars de sécurité de $270^\circ$ sur deux coins opposés (vue
fusionnée à $360^\circ$) ou un lidar en hauteur sur un mât (impossible ici : le bras et le bac
seraient dans son champ). Un second lidar n'a pas été ajouté : coût de simulation (deux capteurs de
vision rendus à chaque pas), risque de collision avec le bras et le bac, fusion des nuages de
points. L'arrière est couvert par les ultrasons et l'arrêt d'urgence ; limite acceptée : une
personne arrivant par l'arrière est détectée tard et ne peut pas être suivie (un ultrason donne une
distance, pas une position ni une vitesse).

\begin{figure}[h]
  \centering
  \figimg{0.45}{couverture_capteurs_atelier_cablage.png}
  \caption{Couverture des capteurs, robot au dock (avant vers le haut) : champ du lidar en bleu,
  cônes d'ultrasons (gris : rien, orange : écho, arc à la distance mesurée), points lidar en rouge.
  L'ultrason arrière entend un opérateur que le lidar ne peut pas voir.}
  \label{fig:couverture}
\end{figure}

\paragraph{Obstacles « imprévus ».} Un point détecté est \emph{expliqué par la carte} s'il tombe à
moins de 15\,cm d'un objet cartographié (carte dilatée de 15\,cm). Les autres points sont des
obstacles imprévus (personnes, objets égarés). Ils sont regroupés dans des cellules de 10\,cm et
\textbf{mémorisés 1,5\,s}, de sorte qu'une personne sortie du champ du lidar (sur le côté du robot)
n'est pas oubliée. Les ultrasons complètent le lidar dans ses angles morts. Leur mesure étant une
simple distance dans un cône large, leur tolérance « expliqué par la carte » est plus grande
(30\,cm).

% ===========================================================================
\section{Carte d'occupation}

\textbf{Méthode :} grille de cellules de 5\,cm ($420\times260$ cellules pour l'atelier) générée à
partir de \code{layout.py}. Un objet occupe une cellule s'il se trouve entre 2\,cm et 0,5\,m de
hauteur (la hauteur du robot) : un écran en hauteur ne bloque pas, son pied si. Les personnes ne
sont \emph{pas} sur la carte : elles bougent, c'est le lidar qui les gère.

Trois couches :
\begin{enumerate}
  \item \textbf{occupée} : les objets ;
  \item \textbf{interdite} (\emph{lethal}) : les objets dilatés par le rayon du robot
        $r = 0{,}40$\,m (dilatation morphologique par un disque, OpenCV). Le \emph{centre} du robot
        doit rester hors de cette zone, ce qui permet de planifier pour un point au lieu d'un
        rectangle ;
  \item \textbf{coût de confort} : à partir de la transformée de distance $d$ (distance de chaque
        cellule à l'obstacle le plus proche),
        \begin{equation}
          c = \operatorname{clip}\!\left(\frac{r + C - d}{C},\,0,\,1\right),\qquad C = 0{,}45\ \text{m},
        \end{equation}
        qui décroît de 1 (contre la zone interdite) à 0 (à 0,85\,m de l'obstacle). Les chemins
        restent ainsi au milieu des allées.
\end{enumerate}

\textbf{Justification :} la grille est la représentation la plus simple pour une carte 2D connue ;
elle se combine directement avec A* et avec les détections du lidar, que l'on « peint » dans une
copie de la carte pour re-planifier.

\begin{figure}[h]
  \centering
  \figimg{0.85}{carte_chemins.png}
  \caption{Carte d'occupation (noir : objets, gris : zone interdite, dégradé : coût de confort)
  et chemins planifiés par A*.}
\end{figure}

% ===========================================================================
\section{Planification globale : A*}

\subsection{Principe}
A* explore les cellules depuis le départ en développant toujours la cellule de plus petit
\begin{equation}
  f(n) = g(n) + h(n),
\end{equation}
où $g(n)$ est le coût du meilleur chemin connu du départ à $n$ et $h(n)$ une estimation du coût
restant. Si $h$ ne surestime jamais le coût réel (heuristique \emph{admissible}), le premier chemin
qui atteint le but est optimal.

\subsection{Choix faits dans ce projet}
\begin{itemize}
  \item \textbf{Grille 8-connexe} : déplacements droits (coût 1) et diagonaux (coût $\sqrt2$).
        On obtient des chemins plus courts et plus naturels qu'en 4-connexité.
  \item \textbf{Heuristique octile}, exacte en l'absence d'obstacles sur une grille 8-connexe
        (donc admissible et très informée) :
        \begin{equation}
          h = (\Delta x + \Delta y) + (\sqrt2 - 2)\,\min(\Delta x, \Delta y).
        \end{equation}
  \item \textbf{Coût de pas pondéré} par le confort : un pas de longueur $\ell$ vers une cellule de
        coût $c$ coûte
        \begin{equation}
          \ell\,(1 + 4\,c).
        \end{equation}
        On obtient un compromis entre chemin court et marge vis-à-vis des obstacles.
        L'heuristique reste admissible, car le coût de pas est toujours $\geq \ell$.
  \item \textbf{Pas de coupe de coin} : un pas diagonal est refusé si l'une des deux cellules
        voisines est interdite.
  \item \textbf{Départ/but dans la zone interdite} : on prend la cellule libre la plus proche.
  \item \textbf{Lissage par ligne de vue} : le chemin brut en escalier est simplifié. Depuis chaque
        point conservé, on saute au point le plus lointain visible en ligne droite sans traverser de
        cellule interdite ni de cellule de coût $> 0{,}6$. Il reste quelques points de passage
        reliés par des segments droits (par exemple 3 points pour Dock $\rightarrow$ Sertissage,
        15,8\,m).
\end{itemize}

\begin{algorithm}[h]
\caption{A* sur la grille (avec coût de confort)}
\begin{algorithmic}[1]
\Require grille (interdite, coût), départ $s$, but $b$
\Ensure liste de cellules de $s$ à $b$, ou échec
\State $g[s] \gets 0$ ;\ \ \textsc{Ouverts} $\gets \{(h(s), s)\}$
\While{\textsc{Ouverts} non vide}
  \State $n \gets$ cellule de plus petit $f$ dans \textsc{Ouverts} (tas binaire)
  \If{$n = b$} \Return chemin reconstruit par les parents \EndIf
  \State marquer $n$ comme fermée
  \For{chaque voisin $m$ de $n$ (8 directions), libre, non fermé, sans coupe de coin}
     \State $g' \gets g[n] + \ell(n,m)\,(1 + 4\,c[m])$
     \If{$g' < g[m]$}
        \State $g[m] \gets g'$ ;\ \ parent$[m] \gets n$ ;\ \ ajouter $(g' + h(m), m)$ à \textsc{Ouverts}
     \EndIf
  \EndFor
\EndWhile
\State \Return échec
\end{algorithmic}
\end{algorithm}

\subsection{Résultats}
Les 8 poses de poste sont libres sur la carte ; les 28 paires de postes ont toutes un chemin ;
le calcul prend environ 15\,ms par trajet. Cette rapidité permet de re-planifier en cours de route
sans pénaliser le robot.

% ===========================================================================
\section{Suivi du chemin et évitement des obstacles}

\subsection{Suivi : point de visée (type « pure pursuit »)}
Le robot vise un point situé $0{,}6$\,m plus loin sur le chemin (la « carotte ») et s'oriente dans
la direction du chemin, pour que le lidar regarde là où le robot va. La vitesse de croisière est de
$0{,}45$\,m/s, modulée par :
\begin{equation}
v = V_{max}\cdot\operatorname{clip}\!\left(\frac{d_{bloc} - d_{stop}}{d_{lent} - d_{stop}},0,1\right)
\cdot \max\!\big(0{,}25,\cos e_\theta\big)\cdot\min\!\left(1,\ \frac{d_{restant}}{0{,}8} + 0{,}3\right),
\end{equation}
avec $d_{bloc}$ la distance (le long du chemin) au premier obstacle dans un couloir de $\pm0{,}5$\,m,
$d_{lent} = 1{,}8$\,m et $d_{stop} = 0{,}75$\,m. On obtient un ralentissement linéaire à l'approche
d'un obstacle, un ralentissement dans les virages et un ralentissement à l'arrivée. Toute détection
à moins de $1{,}2$\,m limite la vitesse à $0{,}2$\,m/s.

\subsection{Machine à états}
\begin{table}[h]
\centering\small
\begin{tabularx}{\linewidth}{@{}lL@{}}
\toprule
\textbf{État} & \textbf{Comportement et transitions} \\
\midrule
\textsc{Suivre} & Suit le chemin. Obstacle dans le couloir à moins de 0,75\,m, ou n'importe
  où à moins de 0,55\,m $\rightarrow$ \textsc{Attendre}. Ultrason dans la direction du mouvement
  $< 0{,}22$\,m $\rightarrow$ arrêt d'urgence. À moins de 0,35\,m du but $\rightarrow$ \textsc{Final}. \\
\textsc{Attendre} & Arrêté. Voie libre $\rightarrow$ \textsc{Suivre}. Bloqué plus de 3\,s
  $\rightarrow$ re-planification A* en ajoutant les obstacles vus (disques de 0,35\,m) à la carte.
  Obstacle qui se rapproche (plus de 10\,cm en 0,5\,s) $\rightarrow$ \textsc{S'écarter}. \\
\textsc{S'écarter} & Se déplace (sans tourner, grâce au mécanum) dans une direction libre sur la
  carte, à l'opposé des obstacles proches, jusqu'à 1\,m de distance $\rightarrow$ \textsc{Dégager}. \\
\textsc{Dégager} & Attend que plus personne ne soit à moins de 1,2\,m pendant 1\,s, puis
  re-planifie depuis la position atteinte. \\
\textsc{Final} & Approche précise de la pose finale (position + cap) avec la loi P. \\
\bottomrule
\end{tabularx}
\caption{États du navigateur.}
\end{table}

\paragraph{Direction d'écartement.} Chaque obstacle proche $p_i$ « repousse » le robot (position
$p$) avec un poids qui croît quand il est proche :
\begin{equation}
\vec a = \sum_i \frac{p - p_i}{\max(\lVert p - p_i\rVert,\,0{,}1)^2}.
\end{equation}
On essaie la direction de $\vec a$ puis des directions décalées de $\pm30^\circ$, $\pm60^\circ$\dots\
jusqu'à en trouver une libre sur la carte et qui éloigne réellement des obstacles.

\textbf{Justification :} une machine à états explicite est prévisible, facile à tester et à
expliquer, et elle suit la stratégie choisie dans le cahier des charges (ralentir, attendre,
re-planifier). Des méthodes plus avancées (fenêtre dynamique DWA, champs de potentiel continus)
restent des évolutions possibles. La prédiction du mouvement des personnes, d'abord prévue comme
évolution, a été ajoutée (section~\ref{sec:suivi}).

% ===========================================================================
\section{Suivi des opérateurs et prédiction de leur mouvement}\label{sec:suivi}

La perception de la section précédente ne voit que des \emph{points} : elle sait qu'il y a
« quelque chose », pas s'il s'agit d'un carton posé ou d'une personne qui marche, ni dans quelle
direction elle va. Le suivi transforme ces points en \emph{objets} avec une position et une vitesse.

\subsection{Regroupement des points}
Les points lidar non expliqués par la carte qui sont à moins de 0,6\,m les uns des autres forment
un même objet (regroupement « single-linkage », par propagation d'étiquettes sur la matrice des
distances). Le centre de l'objet est la moyenne de ses points.

\textbf{Problème rencontré :} le lidar, à 13\,cm du sol, voit les \emph{jambes} des opérateurs. En
pleine foulée les pieds sont à $\approx 0{,}5$\,m l'un de l'autre ; avec un seuil de 0,35\,m, chaque
jambe devenait un objet (l'une presque immobile, l'autre à deux fois la vitesse de marche), d'où
une erreur de vitesse de 0,4\,m/s. Le seuil de 0,6\,m réunit les deux jambes.

\subsection{Filtre de Kalman à vitesse constante}
Chaque objet a son propre filtre (OpenCV \code{cv2.KalmanFilter}). L'état est
$\mathbf{x} = (x, y, v_x, v_y)^\top$, la mesure $\mathbf{z} = (x, y)^\top$ est le centre du groupe :
\begin{equation}
F = \begin{pmatrix} 1&0&\Delta t&0\\ 0&1&0&\Delta t\\ 0&0&1&0\\ 0&0&0&1 \end{pmatrix},\quad
H = \begin{pmatrix} 1&0&0&0\\ 0&1&0&0 \end{pmatrix},\quad
Q = G\,G^\top\sigma_a^2,\ \ G = \begin{pmatrix} \Delta t^2/2&0\\ 0&\Delta t^2/2\\ \Delta t&0\\ 0&\Delta t \end{pmatrix},\quad
R = \sigma_m^2\,I,
\end{equation}
avec $\sigma_a = 1$\,m/s$^2$ (une personne peut changer de vitesse) et $\sigma_m = 0{,}12$\,m (le
centre des jambes bouge pendant la marche). À chaque pas de 50\,ms :
\begin{align}
\text{prédiction : }& \hat{\mathbf{x}}^- = F\hat{\mathbf{x}},\qquad P^- = F P F^\top + Q,\\
\text{correction : }& K = P^- H^\top (H P^- H^\top + R)^{-1},\quad
  \hat{\mathbf{x}} = \hat{\mathbf{x}}^- + K(\mathbf{z} - H\hat{\mathbf{x}}^-),\quad P = (I - K H)P^-.
\end{align}
\textbf{Association :} les paires (objet suivi, groupe détecté) sont traitées de la plus proche à
la plus éloignée (plus proche voisin), à moins de 0,8\,m de la prédiction. Un groupe non associé
crée un nouvel objet ; un objet non revu pendant 1,5\,s est supprimé (entre-temps sa position
continue d'être \emph{prédite} : une personne cachée un instant n'est pas oubliée) ; deux objets à
moins de 0,5\,m l'un de l'autre sont fusionnés.

\textbf{États :} un objet est confirmé après 4 détections ; il est \emph{en mouvement} si sa vitesse
dépasse 0,15\,m/s, \emph{fixe} s'il reste sous ce seuil pendant 2\,s. Il est aussi noté comme
« ayant bougé » s'il s'est éloigné de plus de 0,6\,m de sa première position confirmée : une
personne arrêtée va repartir, un carton non. Ce critère porte sur la distance et non sur la
vitesse : vu de loin puis sous d'autres angles pendant que le robot roule, un carton montre
d'autres faces et le centre de ses points se décale ($\leq 0{,}3$\,m), ce qui ressemblait à une
vitesse. Position prévue dans $t$ secondes : $\mathbf{p}(t) = \mathbf{p} + \mathbf{v}\,t$.

\textbf{Validation} (\code{test\_tracking} : robot immobile 60\,s, comparaison avec les vraies
positions et vitesses lues dans le simulateur ; un opérateur ne compte que s'il est réellement
touché par le lidar) :
\begin{table}[h]
\centering\small
\begin{tabularx}{\linewidth}{@{}lLccc@{}}
\toprule
\textbf{Scène} & \textbf{Opérateur} & \textbf{Suivi} & \textbf{Erreur position} & \textbf{Erreur vitesse} \\
\midrule
Atelier & traverse à l'ouest & 98,5\,\% & 14\,cm & 0,09\,m/s \\
Atelier & traverse à l'est & 96,4\,\% & 13\,cm & 0,08\,m/s \\
Atelier & entre dans l'allée et s'arrête & 100\,\% & 4\,cm & 0,02\,m/s \\
Arène de démonstration & traverse & 100\,\% & 8\,cm & 0,02\,m/s \\
\bottomrule
\end{tabularx}
\caption{Précision du suivi (l'erreur de position est surtout un décalage constant : le lidar voit
l'avant des jambes, le simulateur donne le centre du corps).}
\end{table}

\begin{figure}[h]
  \centering
  \figimg{0.8}{live_atelier_carte.png}
  \caption{Carte temps réel : route A* (bleu), trajet (vert), opérateur suivi (orange) avec une
  flèche vers sa position prévue 1,5\,s plus tard, objet fixe (gris).}
\end{figure}

\subsection{Utilisation dans la navigation}
\begin{itemize}
  \item \textbf{Croisement prévu :} les positions prévues d'une personne sur 3\,s (tous les
        0,25\,s) qui tombent dans le couloir du robot sont un conflit \emph{seulement si} elle y
        arrive au moment où le robot y passe. Le robot atteint un point situé à $d$ le long du
        chemin en $d / v_{max}$ et l'a dépassé à $(d + 0{,}8)/v_{max}$ ; si la personne n'y est
        prévue que plus d'1\,s après, le robot continue. En cas de conflit, il s'arrête 0,5\,m plus
        tôt que pour un obstacle réel, pour attendre en dehors du passage.
  \item \textbf{Déjà dans le passage :} si le robot est déjà sur la trajectoire d'une personne
        encore loin (au moins 2\,s), il ne s'arrête pas sur place : il continue pour en sortir.
  \item \textbf{Personne qui passe :} une personne dont la trajectoire prévue reste à plus de
        0,8\,m du robot ne déclenche ni manœuvre d'écart ni limitation à 0,2\,m/s.
  \item \textbf{Personne qui s'éloigne} sur le chemin : le robot la suit à $\approx 1{,}4$\,m au
        lieu de s'arrêter (comme un régulateur de distance).
  \item \textbf{Objet fixe qui n'a jamais bougé} (un carton) : re-planification après 0,5\,s au
        lieu de 3\,s ; une personne arrêtée, elle, est attendue.
  \item La traînée de points mémorisés derrière une personne qui marche est oubliée (le suivi
        sait où elle est) ; les règles de sécurité sur ce qui est \emph{réellement} là (arrêt à
        0,55\,m devant le robot, arrêt d'urgence par ultrasons) sont inchangées.
\end{itemize}

\subsection{Résultats : avec et sans prédiction}
Les opérateurs de l'atelier ont été disposés pour cette étape : deux d'entre eux \emph{traversent}
l'allée sur les trajets du robot, sans attendre le robot à distance (ils ne s'arrêtent que s'il est
juste devant eux) et avec des pauses de 6\,s aux extrémités ; un troisième entre dans l'allée et y
reste 12\,s. Mission complète (Dock $\rightarrow$ Magasin, chargement $\rightarrow$ Coupe,
déchargement $\rightarrow$ Dock) :

\begin{table}[h]
\centering\small
\begin{tabularx}{\linewidth}{@{}lLL@{}}
\toprule
& \textbf{Sans prédiction} & \textbf{Avec prédiction (réglage final, 2 essais)} \\
\midrule
Durée de la mission & 167,6\,s & 162,4 et 175,4\,s \\
Distance minimale à un opérateur & 0,74--0,77\,m & 0,78--0,86\,m \\
Contacts & 0 & 0 \\
\bottomrule
\end{tabularx}
\caption{Effet de la prédiction avec des opérateurs qui s'écartent au dernier moment.}
\end{table}

Avec des opérateurs qui \emph{ne réagissent jamais} au robot (essai préliminaire), la prédiction
divisait les contacts par trois et faisait passer la distance minimale de 0,07\,m à 0,33--0,46\,m,
sans pouvoir les supprimer tous : le youBot (0,45\,m/s) est plus lent que les opérateurs
(0,7--0,8\,m/s). Le premier réglage de la prédiction s'arrêtait pour tout croisement prévu, même
quand le robot serait passé à temps, et re-planifiait autour des opérateurs en pause : il était
40\,s plus lent ; les deux règles ci-dessus (conflit dans le temps, objet « qui n'a jamais bougé »)
ont corrigé cela. \textbf{Bilan :} la prédiction apporte surtout de la sécurité face à des
personnes inattentives ; bien réglée, elle ne coûte pas de temps et garde plus de marge.

\paragraph{Obstacle fixe absent de la carte.} Un carton est posé dans l'allée sur le chemin du
Sertissage (hors de la carte). Le suivi le classe fixe 2,7\,s après l'avoir vu ; le robot s'arrête
à 0,74\,m et re-planifie 0,5\,s plus tard (détour de 6,8\,m), 0 contact. Il a été placé à l'entrée
de la partie de l'allée menant au Sertissage, car plus près il aurait obligé le robot à attendre
ou à contourner dans le passage de l'opérateur qui traverse.

% ===========================================================================
\section{Accostage de précision : ArUco + lidar}\label{sec:dock}

Après la navigation, le robot n'est connu qu'à environ 2\,cm et $1^\circ$ près, alors que le bras a besoin
d'une précision de l'ordre du millimètre. L'accostage se fait en deux phases, chacune avec le
capteur le mieux adapté.

\subsection{Phase 1 : alignement par la caméra (marqueur ArUco)}
Chaque table de transfert porte un marqueur ArUco (dictionnaire \code{4X4\_50}, 15\,cm, sur fond
blanc). OpenCV détecte ses 4 coins avec un affinage sous-pixel, puis \code{solvePnP} (méthode
\code{IPPE\_SQUARE}, spécifique aux cibles carrées planes) donne sa position et son orientation
dans le repère de la caméra. Les paramètres intrinsèques viennent du simulateur :
$f = \frac{W/2}{\tan(\alpha/2)}$, avec $W = 640$\,px et $\alpha = 60^\circ$. La pose est ramenée dans
le repère du robot à partir de la pose de la caméra lue dans la scène, ce qui intègre automatiquement
l'inclinaison de $15^\circ$.

Le robot se place à $0{,}63$\,m face au marqueur, à l'aide d'un lissage exponentiel des mesures
($x \leftarrow x + 0{,}4\,(x_{mes} - x)$) et d'une loi P. L'identifiant lu \emph{vérifie} en même
temps qu'il s'agit du bon poste. Test : 24/24 détections, erreur de 1\,mm et $2{,}4^\circ$.

\subsection{Phase 2 : approche par le lidar}
De près, la caméra ne voit plus le marqueur en entier. La face avant de la table est une droite
dans le nuage lidar : un ajustement par moindres carrés $f = a + b\,l$ sur les points situés
devant le robot ($|l| < 0{,}25$\,m) donne la distance $a$ et l'angle $-\arctan b$ de la face. Le
robot avance jusqu'à $0{,}40$\,m de la face. Test sur les 7 tables, départ avec $\pm10$\,cm et
$\pm8^\circ$ d'erreur : erreur finale $\leq 5$\,mm et $\leq 0{,}2^\circ$.

\textbf{Justification :} c'est la fusion de deux capteurs complémentaires. La caméra identifie le
poste et corrige les grandes erreurs latérales ; le lidar est très précis en distance et en angle
dans les derniers centimètres.

% ===========================================================================
\section{Manipulation : bras du youBot}

\subsection{Cinématique inverse analytique}
Les articulations 1 à 3 forment un bras plan à 3 segments ($L_1 = 0{,}155$, $L_2 = 0{,}135$,
$L_3 = 0{,}19$\,m jusqu'au point de prise) ; l'articulation 0 oriente ce plan. Pour un point de
prise $P$ et une inclinaison de pince $\varphi$ :
\begin{enumerate}
  \item angle de base : $j_0 = \operatorname{atan2}(\Delta y, \Delta x)$ ;
  \item poignet : $W = P - L_3(\cos\varphi, -\sin\varphi)$ dans le plan du bras ;
  \item triangle épaule--coude--poignet de côtés $L_1, L_2, d = \lVert SW\rVert$ : loi des cosinus,
        \begin{equation}
           \gamma = \arccos\frac{L_1^2 + d^2 - L_2^2}{2\,L_1\,d},
        \end{equation}
        solution « coude en haut » ;
  \item on essaie d'abord l'inclinaison la plus forte, puis on la réduit par pas de $2^\circ$
        jusqu'à une solution qui respecte les butées articulaires.
\end{enumerate}
La cinématique directe sert à vérifier la cinématique inverse (test : erreur $\leq 5$\,mm).
\textbf{Justification :} pour un bras plan, la solution analytique est exacte, instantanée et sans
problème de convergence, contrairement à un solveur numérique.

\subsection{Prise, transport et dépose}
Les bacs ont une poignée (petit poteau de $3\times3\times5$\,cm) que la pince serre. Les
mouvements articulaires sont interpolés (profil « smoothstep » $3s^2 - 2s^3$, sans à-coups). La
prise est confirmée par l'ouverture des doigts (ils s'arrêtent sur la poignée). Le bac est alors
\emph{verrouillé} à la main pendant le transport, ce qui évite qu'il pivote autour de la poignée,
puis relâché à la dépose avec la même inclinaison qu'à la prise, de façon à ce qu'il se pose bien à
plat. Le bac est porté sur la plateforme arrière du robot (base du bras tournée à $167^\circ$).

% ===========================================================================
\section{Cellule de contrôle qualité}

\begin{itemize}
  \item \textbf{Décision par vision} : une caméra au-dessus du convoyeur d'inspection lit
        l'étiquette du bac. Segmentation par seuils HSV (vert = conforme, rouge = défaut, le rouge
        étant à cheval sur $H = 0/180$), puis composantes connexes (aire minimale de 40\,px). La
        décision est prise quand l'étiquette passe au milieu de l'image, puis rangée dans une
        \textbf{file FIFO} : les bacs ne peuvent pas se doubler sur un convoyeur.
  \item \textbf{Bras UR5 + ventouse} : les mouvements courts (descente vers un bac, remontée) sont
        des lignes droites cartésiennes. À chaque pas, la cible avance un peu sur la ligne et simIK
        résout la cinématique inverse par \emph{moindres carrés amortis} :
        \begin{equation}
          \Delta q = J^\top\big(J J^\top + \lambda^2 I\big)^{-1} e,
        \end{equation}
        stable près des singularités. Le calcul se fait dans un monde IK séparé, plus robuste que
        la synchronisation avec le bras dynamique.
  \item \textbf{Transferts convoyeur $\leftrightarrow$ tables dans l'espace articulaire} : sur une
        ligne droite passant près de la base de l'UR5, l'IK repliait le bras et son avant-bras
        venait appuyer sur le convoyeur ; l'articulation de base restait bloquée et le bac était
        relâché au-dessus du tapis. Désormais l'IK calcule seulement les angles \emph{finaux} et
        toutes les articulations tournent ensemble de leur angle de départ à leur angle d'arrivée.
        La somme épaule + coude + poignet variant linéairement, la ventouse reste orientée vers le
        bas. La hauteur de passage au-dessus des tables a été abaissée (à 0,68\,m, la position était
        à la limite de portée du bras).
  \item \textbf{Vérification de la position réelle} : le monde IK atteint toujours sa cible, pas
        forcément le bras réel. Après chaque mouvement avec un bac, la position réelle de la
        ventouse est comparée à la cible ; au-delà de 3\,cm, la cellule s'arrête en défaut au lieu
        de relâcher le bac ailleurs (comme une cellule industrielle).
  \item \textbf{Une place par table} : la seconde place, derrière la première et trop près de la
        base du bras (0,55\,m), obligeait le bras à se replier sur son propre socle. Chaque table a
        donc une seule place ; si elle est occupée, le bac suivant attend en bout de convoyeur que
        l'AMR ait emporté le premier.
  \item \textbf{Arrêt d'accumulation} : une barrière lumineuse en bout de convoyeur arrête le tapis
        quand un bac arrive et le relance quand le bras l'a soulevé. Les échecs de prise par la
        ventouse sont passés de 21 à 0.
  \item \textbf{Exécution en parallèle} : la cellule est écrite comme un générateur Python
        (co-routine). Elle avance d'un pas avant chaque pas de simulation, quoi que fasse le robot
        (rouler, s'accoster, manipuler), comme un automate indépendant.
\end{itemize}

% ===========================================================================
\section{Validation et résultats}

\begin{table}[h]
\centering\small
\begin{tabularx}{\linewidth}{@{}lL@{}}
\toprule
\textbf{Test} & \textbf{Résultat} \\
\midrule
Déplacement mécanum & 0,4\,m $\rightarrow$ 0,392 / 0,388\,m ; 1,0\,rad $\rightarrow$ 1,006\,rad \\
Carré de 1\,m & erreur réelle de 2 à 6\,cm (limitée par le bruit UWB) \\
Capteurs & lidar 0,950 / 0,641\,m (attendu 0,95 / 0,64) ; ultrasons 0,65 / 0,42 / 0,65\,m \\
Planification & 28/28 trajets, environ 15\,ms par trajet \\
Navigation avec opérateurs & première disposition : 3 trajets, 0 contact, distance minimale 0,73\,m \\
Suivi des opérateurs & 96--100\,\% du temps où ils sont visibles, erreur de vitesse 0,02--0,09\,m/s \\
Prédiction & 0 contact, distance minimale 0,78--0,86\,m (0,74--0,77\,m sans prédiction), même durée \\
Carton absent de la carte & classé fixe, re-planification 0,5\,s après l'arrêt, 0 contact \\
ArUco & 24/24 détections, 1\,mm / $2{,}4^\circ$ \\
Accostage & 7 tables, $\leq 5$\,mm / $\leq 0{,}2^\circ$ \\
Prise / dépose & bac posé droit sur la table cible \\
Mission complète & Dock $\rightarrow$ Magasin (chargement) $\rightarrow$ Coupe (déchargement)
  $\rightarrow$ Dock : 162--176\,s simulées, 0 contact, deux opérateurs traversant l'allée ; en
  parallèle, cellule CQ : un bac sur chaque table, le second bac conforme attend en bout de
  convoyeur \\
Arène de démonstration & 9/9 étapes réussies, 186\,s simulées, 0 contact \\
\bottomrule
\end{tabularx}
\caption{Synthèse des tests.}
\end{table}

\begin{figure}[h]
  \centering
  \figimg{0.6}{demo_vue_dessus.png}
  \caption{Arène de démonstration du youBot : une zone par capacité (capteurs, mécanum, chicane,
  postes A/B avec marqueurs ArUco, passage d'un opérateur, carton absent de la carte, dock).}
\end{figure}

% ===========================================================================
\section{Limites et perspectives}

\begin{itemize}
  \item \textbf{Personnes :} le robot prédit leur mouvement (vitesse constante) mais reste plus
        lent qu'elles (0,45 contre 0,7--0,8\,m/s) : face à une personne qui ne fait \emph{jamais}
        attention à lui, il peut encore être rattrapé, par exemple quand elle sort de derrière une
        machine à moins d'un mètre de son chemin, ou quand il est déjà dans son passage. Un
        modèle de déplacement plus riche (personnes qui tournent, s'arrêtent) et une vitesse de
        robot plus élevée sont des pistes.
  \item \textbf{Arrière du robot :} sans lidar arrière, une personne qui arrive par derrière n'est
        vue que par les ultrasons, tard et sans être suivie ; un second lidar de coin serait
        l'évolution naturelle.
  \item \textbf{Suivi :} un carton vu pendant que le robot tourne autour peut paraître bouger
        lentement ($\approx 0{,}2$\,m/s) ; cela ne change pas la décision (il n'a jamais parcouru
        0,6\,m) mais apparaît sur l'affichage.
  \item \textbf{Ultrasons :} leur cône large place parfois l'écho devant un mur vu en biais
        (fausse détection). Le robot est alors simplement prudent ; une meilleure modélisation du
        cône réduirait ces ralentissements.
  \item \textbf{Localisation :} la localisation UWB est simulée ; un recalage lidar sur la carte
        (AMCL) rendrait le robot plus autonome vis-à-vis de l'infrastructure.
  \item \textbf{Cellule de contrôle qualité :} une seule place par table ; le bac suivant attend
        en bout de convoyeur. Une table plus éloignée de la base de l'UR5 permettrait deux places.
  \item \textbf{Intégration :} le gestionnaire de missions (ordres, batterie, ordres de reprise
        automatiques depuis la table \textsc{Rebuts}) et l'interface de supervision constituent les
        prochaines étapes ; un passage à ROS~2 (\code{simROS2}) est possible sans changer les
        algorithmes.
\end{itemize}

\end{document}
